\documentclass[11pt,a4paper,twoside,openright]{report}

% Preamble condiviso — FdC Railway Manager
\usepackage[utf8]{inputenc}
\usepackage[T1]{fontenc}
\usepackage[italian]{babel}
\usepackage{lmodern}
\usepackage{microtype}
\usepackage{geometry}
\geometry{a4paper, margin=2.5cm}

\usepackage{graphicx}
\usepackage{booktabs}
\usepackage{longtable}
\usepackage{tabularx}
\usepackage{hyperref}
\usepackage{xcolor}
\usepackage{listings}
\usepackage{amsmath}
\usepackage{tikz}
\usetikzlibrary{arrows.meta, positioning, shapes.geometric}

\hypersetup{
  colorlinks=true,
  linkcolor=blue!60!black,
  urlcolor=blue!70!black,
  pdftitle={FdC Railway Manager — Documentazione moduli},
  pdfauthor={FDC-App}
}

\definecolor{codebg}{RGB}{245,245,245}
\definecolor{codeframe}{RGB}{200,200,200}
\definecolor{swiftkw}{RGB}{170,13,130}

\lstdefinelanguage{Swift}{
  morekeywords={import,struct,enum,class,func,let,var,if,else,guard,
    return,async,await,throws,try,catch,init,public,private,static,
    protocol,extension,true,false,nil},
  sensitive=true,
  morecomment=[l]{//},
  morecomment=[s]{/*}{*/},
  morestring=[b]",
}

\lstset{
  language=Swift,
  basicstyle=\ttfamily\small,
  backgroundcolor=\color{codebg},
  frame=single,
  rulecolor=\color{codeframe},
  breaklines=true,
  columns=fullflexible,
  keepspaces=true,
  showstringspaces=false,
  keywordstyle=\color{swiftkw}\bfseries,
  commentstyle=\itshape\color{gray},
}

\newcommand{\repo}[1]{\texttt{#1}}
\newcommand{\module}[1]{\textbf{#1}}
\newcommand{\wip}{\textcolor{orange}{\textbf{[WIP]}}}


\title{%
  \textbf{FdC Railway Manager}\\[0.4em]
  \large Documentazione modulare
}
\author{Progetto FDC-App}
\date{Luglio 2026 --- v0.1}

\begin{document}

\maketitle
\tableofcontents
\listoffigures
\listoftables

\part{Fondamenti}
\chapter{Prefazione}

\section{Scopo del documento}

Questo manuale descrive i \textbf{moduli software} del progetto
\emph{FdC Railway Manager}: architettura, responsabilità, dipendenze
e flussi dati. È pensato per sviluppatori che contribuiscono al
refactor verso un \emph{functional core} testabile.

La documentazione viene estesa \textbf{incrementalmente}: ogni capitolo
corrisponde a un modulo (o sotto-modulo) e può essere aggiornato in
autonomia senza riscrivere l'intero volume.

\section{Pubblico e prerequisiti}

\begin{itemize}
  \item Conoscenza di Swift 5.9+ e SwiftUI.
  \item Familiarità con pattern \emph{Strangler}, CQS e DI esplicita.
  \item Nozioni base di ferrovie: stazione, binario, Taktfahrplan.
\end{itemize}

\section{Convenzioni}

\begin{description}
  \item[\repo{Domain/Model/}] Tipi di dominio condivisi; preferibilmente
    solo \texttt{Foundation} (o \texttt{CoreLocation} dove necessario).
  \item[\repo{Services/}] Logica applicativa orchestrata dai ViewModel.
  \item[\repo{Models/}] Adapter, \texttt{NetworkModel}, estensioni UI.
  \item[Pipeline] Sequenza deterministica di trasformazioni su \texttt{[Train]}.
\end{description}

\section{Versione}

\begin{center}
\begin{tabular}{ll}
  \toprule
  \textbf{Versione documento} & 0.1 (Fase 4 completata) \\
  \textbf{Ultimo aggiornamento}  & Luglio 2026 \\
  \textbf{Commit di riferimento} & \texttt{00db4c2} \\
  \bottomrule
\end{tabular}
\end{center}

\chapter{Architettura generale}
\label{ch:architettura}

\section{Visione}

L'applicazione è organizzata in tre aree funzionali indipendenti
(\module{Editor}, \module{Scheduler}, \module{Simulator}) che
comunicano esclusivamente tramite tipi in \repo{Domain/Model/}.
Le dipendenze puntano sempre verso il dominio (regola delle frecce).

\section{Struttura repository}

\begin{figure}[htbp]
\centering
\begin{tikzpicture}[
  box/.style={draw, rounded corners, minimum width=3.2cm,
    minimum height=0.9cm, align=center, font=\small},
  arr/.style={-{Stealth[length=2mm]}, thick}
]
  \node[box, fill=blue!8] (ui) {UI / ViewModel};
  \node[box, fill=green!10, below=1.2cm of ui] (svc) {Services};
  \node[box, fill=orange!12, below=1.2cm of svc] (dom) {Domain/Model};
  \node[box, fill=gray!15, right=2.8cm of svc] (infra) {Infrastructure};
  \node[box, fill=gray!15, left=2.8cm of svc] (models) {Models (bridge)};

  \draw[arr] (ui) -- (svc);
  \draw[arr] (svc) -- (dom);
  \draw[arr] (infra) -- (dom);
  \draw[arr] (models) -- (dom);
  \draw[arr, dashed] (ui) -- (models);
\end{tikzpicture}
\caption{Direzione delle dipendenze (semplificata).}
\end{figure}

\section{Refactor Strangler (Fasi 0--4)}

Il monolite \texttt{RailwayScheduleOptimizer} è stato sostituito da
\repo{Services/Scheduling/}. La migrazione dei modelli da
\repo{Models/} a \repo{Domain/Model/} è in corso: i tipi infrastrutturali
(\texttt{Node}, \texttt{Edge}, \texttt{Signal}, \ldots) risiedono già
nel dominio; \texttt{NetworkModel} e i DTO restano nel layer bridge.

\section{Swift Packages (scaffold)}

\begin{table}[htbp]
\centering
\caption{Pacchetti SPM locali (\repo{Packages/}).}
\begin{tabular}{lll}
  \toprule
  \textbf{Package} & \textbf{Stato} & \textbf{Note} \\
  \midrule
  \texttt{FDCDomain}   & Compilabile & 9 tipi Foundation via symlink \\
  \texttt{FDCScheduling} & Scaffold & Non linkato al target Xcode \\
  \bottomrule
\end{tabular}
\end{table}

\section{Test}

\begin{itemize}
  \item \textbf{20 test XCTest} nel target \texttt{FdC Railway ManagerTests}.
  \item \textbf{2 test} nel package \texttt{FDCDomain}.
  \item I servizi di scheduling sono \texttt{struct} pure su
    \texttt{RailwayTopology} (no \texttt{@MainActor}, no singleton).
\end{itemize}

\section{Feature flags}

Funzionalità sperimentali in \repo{FeatureFlags.swift}:
\texttt{simulatorGameMode}, \texttt{realtimeConflictHighlight}.
Il flag \texttt{useLegacyScheduleOptimizer} è stato rimosso con la Fase 4.


\part{Moduli core}
\chapter{Modulo FDCDomain}
\label{ch:fdcdomain}

\section{Responsabilità}

\module{FDCDomain} contiene i \textbf{tipi di dominio ferroviario}:
entità immutabili (o quasi), \texttt{Codable} per persistenza,
senza dipendenze da SwiftUI o da \texttt{NetworkModel}.

Il package SPM in \repo{Packages/FDCDomain/} espone un sottoinsieme
\textbf{Foundation-only} compilabile standalone; il target Xcode include
l'intera cartella \repo{FdC Railway Manager/Domain/Model/}.

\section{Grafo del dominio}

\begin{figure}[htbp]
\centering
\begin{tikzpicture}[
  type/.style={draw, ellipse, minimum width=2.4cm, minimum height=0.8cm,
    align=center, font=\footnotesize},
  arr/.style={-{Stealth[length=1.5mm]}}
]
  \node[type] (topo) {RailwayTopology};
  \node[type, below left=1cm and 0.5cm of topo] (node) {Node};
  \node[type, below right=1cm and 0.5cm of topo] (edge) {Edge};
  \node[type, below=1.8cm of node] (train) {Train};
  \node[type, below=0.6cm of train] (stop) {RelationStop};
  \node[type, right=2.2cm of edge] (seg) {TrackSegment};
  \node[type, above=0.5cm of seg] (sig) {Signal};
  \node[type, below=0.6cm of edge] (tcp) {TrackControlPoint};
  \node[type, left=1.8cm of node] (elec) {ElectrificationType};
  \node[type, right=1.8cm of node] (rc) {RoutingConstraint};

  \draw[arr] (topo) -- (node);
  \draw[arr] (topo) -- (edge);
  \draw[arr] (edge) -- (seg);
  \draw[arr] (edge) -- (tcp);
  \draw[arr] (seg) -- (sig);
  \draw[arr] (node) -- (elec);
  \draw[arr] (node) -- (rc);
  \draw[arr] (train) -- (stop);
\end{tikzpicture}
\caption{Relazioni principali tra tipi di dominio.}
\end{figure}

\section{Tipi per categoria}

\subsection{Topologia}

\begin{description}
  \item[\texttt{RailwayTopology}] Snapshot immutabile di \texttt{nodes}
    e \texttt{edges}. Delega pathfinding a \texttt{NetworkPathfinder}
    (Fase 5). Nessun \texttt{init(network:)}.
  \item[\texttt{Node}] Stazione, deposito, interscambio; vincoli di
    instradamento, altitudine, parametri Takt (\texttt{taktMinutes}).
  \item[\texttt{Edge}] Tratta tra due nodi; segmenti \texttt{TrackSegment},
    punti di controllo, elettrificazione, limite velocità.
\end{description}

\subsection{Infrastruttura di dettaglio}

\begin{description}
  \item[\texttt{TrackSegment}] Binario fisico con ordine, lunghezza,
    segnale opzionale, occupazione.
  \item[\texttt{TrackControlPoint}] Punto sulla tratta (lat/lon/alt).
  \item[\texttt{Signal}] Segnale con aspetto (\texttt{stop}, \texttt{proceed},
    \texttt{caution}).
  \item[\texttt{Switch}] Deviatione con \texttt{nodeId} e collegamenti.
  \item[\texttt{ElectrificationType}] Trazione: diesel, 3\,kV, 25\,kV, 950\,V.
  \item[\texttt{RoutingConstraint}] Vincoli di percorso per stazione.
\end{description}

\subsection{Servizio e rotazione}

\begin{description}
  \item[\texttt{Train}] Treno con fermate \texttt{RelationStop}, orari,
    assegnazione binari, \texttt{routeId}.
  \item[\texttt{RelationStop}] Fermata o transito su una stazione.
  \item[\texttt{Vehicle}] Rotazione materiale rotabile.
  \item[\texttt{TrainRoute}] Template di percorso (origine, destinazione,
    sequenza stazioni). \texttt{displayColor} in \repo{Models/TrainRoute+UI.swift}.
\end{description}

\section{Package SPM}

\begin{lstlisting}
Packages/FDCDomain/
  Package.swift
  Sources/FDCDomain/    # symlink ai .swift Foundation-only
  Tests/FDCDomainTests/
\end{lstlisting}

Tipi inclusi nel package (v0.1):
\texttt{ElectrificationType}, \texttt{Signal}, \texttt{Switch},
\texttt{RoutingConstraint}, \texttt{RelationStop}, \texttt{Vehicle},
\texttt{TrackSegment}, \texttt{TrackControlPoint}, \texttt{RailwayConstants}.

Tipi \textbf{esclusi} fino a decoupling UI / NetworkModel:
\texttt{Node}, \texttt{Edge}, \texttt{Train}, \texttt{RailwayTopology},
\texttt{LineProfilePoint}, \texttt{TrainRoute}, \texttt{Ferrovia}.

\section{Regole di evoluzione}

\begin{enumerate}
  \item Nuovi modelli → \repo{Domain/Model/}, solo \texttt{import Foundation}.
  \item Estensioni SwiftUI → file \texttt{+UI.swift} in \repo{Models/}.
  \item Rimuovere voci duplicate da \texttt{project.pbxproj} (sync group).
  \item Aggiungere test in \texttt{FDCDomainTests} quando il tipo è standalone.
\end{enumerate}

\chapter{Modulo FDCScheduling}
\label{ch:fdcscheduling}

\section{Responsabilità}

\module{FDCScheduling} orchestrate la \textbf{generazione e ottimizzazione
degli orari}. Il codice attivo risiede in
\repo{FdC Railway Manager/Services/Scheduling/}; il package SPM
\repo{Packages/FDCScheduling/} è uno scaffold per l'estrazione futura.

\section{Componenti principali}

\begin{table}[htbp]
\centering
\caption{Servizi di scheduling.}
\begin{tabularx}{\textwidth}{lX}
  \toprule
  \textbf{Tipo} & \textbf{Ruolo} \\
  \midrule
  \texttt{ScheduleGenerationEngine} & Orchestratore: crea treni grezzi,
    invoca pipeline, gestisce modalità Takt/libera. \\
  \texttt{ScheduleOptimizationPipeline} & Pipeline a 8 step (vedi \S\ref{sec:pipeline}). \\
  \texttt{TaktEngine} & Allineamento Taktfahrplan, diagnostica hub. \\
  \texttt{ScheduleRefresher} & Ricalcolo orari fisici (arrivi/partenze). \\
  \texttt{ScheduleAIResolver} & Step AI cloud (\texttt{RailwayAIService}). \\
  \texttt{KinematicCalculator} & Tempi di percorrenza su topologia. \\
  \texttt{PathResolver} & Risoluzione percorsi tra stazioni. \\
  \texttt{VehicleSuitabilityEngine} & Compatibilità rotabile/linea. \\
  \bottomrule
\end{tabularx}
\end{table}

\section{Pipeline a 8 step}
\label{sec:pipeline}

\begin{figure}[htbp]
\centering
\begin{tikzpicture}[
  step/.style={draw, rectangle, rounded corners, minimum width=4.8cm,
    minimum height=0.75cm, align=center, font=\small},
  arr/.style={-{Stealth[length=2mm]}}
]
  \node[step, fill=yellow!15] (s1) {1. Shift partenze (greedy)};
  \node[step, fill=blue!10, below=0.5cm of s1] (s2) {2. Refresh fisico};
  \node[step, fill=green!12, below=0.5cm of s2] (s35) {3--5. TaktEngine};
  \node[step, fill=purple!10, below=0.5cm of s35] (s6) {6. AI cloud (opz.)};
  \node[step, fill=orange!12, below=0.5cm of s6] (s7) {7. Ottimizzatore genetico};
  \node[step, fill=red!8, below=0.5cm of s7] (s8) {8. Verifica finale};

  \draw[arr] (s1) -- (s2);
  \draw[arr] (s2) -- (s35);
  \draw[arr] (s35) -- (s6);
  \draw[arr] (s6) -- (s7);
  \draw[arr] (s7) -- (s8);
\end{tikzpicture}
\caption{\texttt{ScheduleOptimizationPipeline.execute}.}
\end{figure}

\subsection{Dettaglio step}

\begin{enumerate}
  \item \textbf{Ottimizzazione partenze} --- Per ogni nuovo treno, prova
    shift coarse ($\pm 10, 20, 30, 50$\,min) e fine ($\pm 1, 2, 5$\,min);
    minimizza i conflitti rilevati da \texttt{ConflictManager}.
  \item \textbf{Refresh fisico} --- \texttt{ScheduleRefresher.refreshMultiple}
    ricalcola tempi di transito e fermata.
  \item \textbf{Allineamento Takt} --- Se almeno un nodo ha
    \texttt{taktMinutes}, \texttt{TaktEngine.generaOrarioCadenzato}
    allinea arrivi/partenze all'hub; GA disabilitato in modalità Takt.
  \item \textbf{AI cloud} --- Opzionale (\texttt{useAI}); delega a
    \texttt{ScheduleAIResolver} e \texttt{RailwayAIService}.
  \item \textbf{GA} --- \texttt{GeneticOptimizer} con 250 iterazioni
    (solo se non Takt).
  \item \textbf{Verifica} --- Secondo refresh + report conflitti residui.
\end{enumerate}

\section{Entry point}

\begin{lstlisting}
// ScheduleGenerationEngine (semplificato)
let topology = RailwayTopology(nodes: network.nodes, edges: network.edges)
let pipeline = ScheduleOptimizationPipeline(topology: topology)
return await pipeline.execute(
    newTrains: trains,
    existingTrains: existing,
    useAI: false,
    useGA: useGA,
    preferredTaktNodeId: taktNode
)
\end{lstlisting}

\section{Dipendenze esterne (non ancora nel package)}

\begin{itemize}
  \item \texttt{ConflictManager} --- Rilevamento conflitti capacità/binari.
  \item \texttt{GeneticOptimizer} --- Meta-euristica su orari.
  \item \texttt{RailwayAIService} --- Ottimizzazione remota.
  \item \texttt{NetworkModel} --- Bridge verso UI e persistenza.
\end{itemize}

L'estrazione in \texttt{FDCScheduling} richiederà protocolli in
\repo{Domain/Services/} (pianificato Fase 5+).

\section{Test di regressione}

\begin{table}[htbp]
\centering
\caption{Suite test scheduling (\texttt{FdC Railway ManagerTests}).}
\begin{tabular}{ll}
  \toprule
  \textbf{File} & \textbf{Copertura} \\
  \midrule
  \texttt{ServicesSchedulingTests} & Pipeline, refresher, path resolver \\
  \texttt{Taktfahrplan120MinTests} & Cicli Takt 120\,min, hub timing \\
  \texttt{TrackAssignmentTests}   & Assegnazione binari \\
  \texttt{EdgeMigrationTests}     & Migrazione modello Edge \\
  \bottomrule
\end{tabular}
\end{table}

\section{Storico}

Il capitolo sostituisce la documentazione del legacy
\texttt{RailwayScheduleOptimizer} (rimosso in commit \texttt{00db4c2}).
Per confronto algoritmico vedere anche
\repo{FdC Railway Manager/RailwayAlgorithmDocs.md}.

\chapter{Modulo Pathfinding}
\label{ch:pathfinding}

\section{Responsabilità}

Il pathfinding calcola percorsi minimi sulla topologia ferroviaria:
Dijkstra, distanza cumulata, percorsi alternativi, ordinamento nodi per linea
fisica. Dal refactor Fase 5 la logica risiede in
\repo{Domain/Pathfinding/NetworkPathfinder.swift} --- puro Foundation,
senza \texttt{NetworkModel}.

\section{API principale}

\begin{table}[htbp]
\centering
\caption{Funzioni \texttt{NetworkPathfinder}.}
\begin{tabularx}{\textwidth}{lX}
  \toprule
  \textbf{Funzione} & \textbf{Scopo} \\
  \midrule
  \texttt{findShortestPath} & Dijkstra: sequenza nodi + distanza \\
  \texttt{findPathEdges}    & Percorso come lista di \texttt{Edge} \\
  \texttt{calculatePathDistance} & Somma \texttt{edge.distance} su un path \\
  \texttt{findAlternativePaths}  & Varianti Rapido / Alternativo / Panoramico \\
  \texttt{findLineEndpoints}     & Capolinea (grado 1 nel sottografo) \\
  \texttt{buildOrderedPath}        & Ordinamento topologico nodi linea \\
  \texttt{resolveLinePath}         & Validazione path linea fisica \\
  \bottomrule
\end{tabularx}
\end{table}

\section{Integrazione}

\begin{itemize}
  \item \texttt{RailwayTopology} delega \texttt{findPathEdges} e
    \texttt{calculatePathDistance} a \texttt{NetworkPathfinder}.
  \item \texttt{NetworkModel+Pathfinding} resta come \emph{facciata} per il
    layer UI (Strangler): inoltra tutte le chiamate al dominio.
  \item Rimosso \texttt{RailwayTopology.init(network:)} --- nessuna
    dipendenza Domain $\rightarrow$ \texttt{NetworkModel}.
\end{itemize}

\section{Vincoli algoritmici}

\begin{itemize}
  \item Binario singolo non accoppiato: tratto \textbf{non direzionale}.
  \item Coppie direzionali (\texttt{pairedEdgeId}): rispetto del verso.
  \item \texttt{restrictIntermediateStations}: Dijkstra non espande \emph{attraverso}
    stazioni intermedie (solo origine).
\end{itemize}

\section{Test}

\texttt{PathfindingTests} (target app): percorso diretto A$\rightarrow$B e
delega da \texttt{RailwayTopology}.

\section{Evoluzione (Fase 6+)}

Inclusione in \texttt{Packages/FDCDomain} quando \texttt{Node} sarà
decoupled da \texttt{.localized}. Pathfinding è già Foundation-only.

\chapter{Integrazione Swift Package}
\label{ch:spm}

\section{FDCDomain nel target Xcode}

Dal refactor Fase 6 il package locale \texttt{Packages/FDCDomain}
è collegato al target \emph{FdC Railway Manager}:

\begin{itemize}
  \item \texttt{XCLocalSwiftPackageReference} con \texttt{relativePath}
  \item \texttt{packageProductDependencies} + fase \texttt{Frameworks}
  \item \texttt{PBXFileSystemSynchronizedBuildFileExceptionSet}: i sorgenti
    in \repo{Domain/Model/} e \repo{Domain/Pathfinding/} non vengono
    ricompilati nell'app (evita simboli duplicati).
  \item \repo{Models/DomainBridge.swift} con \texttt{@\_exported import FDCDomain}
    e typealias per compatibilità con il codice esistente.
\end{itemize}

\section{Decoupling UI (Fase 6)}

\begin{table}[htbp]
\centering
\caption{Spostamenti layer UI.}
\begin{tabularx}{\textwidth}{lX}
  \toprule
  \textbf{Prima} & \textbf{Dopo} \\
  \midrule
  \texttt{Node.localizedName} & \texttt{Models/Node+UI.swift} \\
  \texttt{TrainDatabaseManager + View} & \texttt{Infrastructure/} + \texttt{UI/Views/TrainDatabase/} \\
  \texttt{TrainDatabaseModels} (SwiftUI) & Solo struct \texttt{Codable} in Domain \\
  \bottomrule
\end{tabularx}
\end{table}

\section{API pubbliche}

I tipi esposti dal package richiedono \texttt{public} su proprietà e metodi
usati dall'app, ad esempio:

\begin{itemize}
  \item \texttt{RailwayTopology.nodes/edges}
  \item \texttt{Node.getTracksByProvenance}, \texttt{defaultVisualType}
  \item \texttt{Train.getPreferredTracks}
  \item \texttt{NetworkPathfinder} (tutti gli entry point statici)
\end{itemize}

\section{Test}

\begin{itemize}
  \item \textbf{5 test} in \texttt{FDCDomainTests} (package SPM)
  \item \textbf{24 test} nel target app (inclusi \texttt{PathfindingTests})
\end{itemize}

\section{Prossimi passi (Fase 8)}

\begin{itemize}
  \item Estrarre \texttt{FDCScheduling} con protocolli in \texttt{Domain/Services/}
  \item Capitolo LaTeX \texttt{06-network-adapter.tex} completato in Fase 7
\end{itemize}

\chapter{NetworkModel come adapter UI}
\label{ch:network-adapter}

\section{Obiettivo (Fase 7)}

Separare la \textbf{shell osservabile} usata dall'interfaccia dalla logica
di query sulla topologia. \texttt{NetworkModel} resta nel target app;
le operazioni pure vivono in \module{FDCDomain}.

\section{NetworkTopologyService}

\texttt{NetworkTopologyService} (\repo{Domain/Services/}) è una struct
\texttt{Sendable} che incapsula \texttt{RailwayTopology} e offre:

\begin{itemize}
  \item lookup: \texttt{findNode}, \texttt{findEdge}
  \item filtri: \texttt{allStations}, \texttt{allJunctions}
  \item ordinamento: \texttt{sortedNodes}, \texttt{sortedEdges}
  \item connettività: \texttt{getConnectedNodeIds}, \texttt{getNeighborStations}
  \item geometria: \texttt{calculateDistance} (hub-aware, km)
\end{itemize}

Nessuna mutazione, nessuna dipendenza da SwiftUI o undo.

\section{NetworkModel}

\texttt{NetworkModel} delega le query a \texttt{NetworkTopologyService}
tramite una computed property \texttt{topologyService}. Espone
\texttt{topology: RailwayTopology} per i servizi di scheduling.

Le mutazioni (aggiunta nodi, binari, undo) restano in
\repo{Models/Network/NetworkModel+Updates.swift}.

\section{DomainBridge}

\texttt{DomainBridge.swift} espone il package e risolve l'ambiguità
\texttt{Edge} (SwiftUI vs dominio):

\begin{verbatim}
@_exported import FDCDomain
public typealias Edge = FDCDomain.Edge
\end{verbatim}

I typealias legacy per gli altri tipi sono stati rimossi.

\section{TrainCategory e DTO}

\begin{itemize}
  \item \texttt{TrainCategory} --- enum Foundation in \repo{Domain/Model/};
        estensione SwiftUI in \repo{Models/TrainCategory+UI.swift}
  \item \texttt{RailwayNetworkDTO} --- struct \texttt{Codable} in
        \repo{Infrastructure/} (I/O JSON, non dominio puro)
\end{itemize}

\section{Test}

\begin{itemize}
  \item \textbf{2 test} aggiuntivi in \texttt{PathfindingTests} (app)
  \item \textbf{2 test} in \texttt{FDCDomainTests} per \texttt{NetworkTopologyService}
\end{itemize}

\section{Prossimi passi (Fase 8)}

\begin{itemize}
  \item Estrarre \texttt{FDCScheduling} in package SPM
  \item Protocolli in \texttt{Domain/Services/} per ConflictManager, GA, AI
\end{itemize}

\chapter{Package FDCScheduling}
\label{ch:scheduling-package}

\section{Obiettivo (Fase 8)}

Estrarre la pipeline di scheduling in \module{FDCScheduling} SPM,
collegata a \module{FDCDomain} tramite protocolli in \texttt{Domain/Services/}.

\section{Architettura dipendenze}

\begin{figure}[htbp]
\centering
\begin{tikzpicture}[
  box/.style={draw, rectangle, rounded corners, minimum width=3.2cm,
    minimum height=0.8cm, align=center, font=\footnotesize},
  arr/.style={-{Stealth[length=1.5mm]}}
]
  \node[box, fill=blue!10] (domain) {FDCDomain};
  \node[box, fill=green!12, right=1.8cm of domain] (sched) {FDCScheduling};
  \node[box, fill=orange!12, right=1.8cm of sched] (app) {App target};

  \node[box, fill=gray!10, below=1.2cm of sched] (proto) {Protocolli\\Conflict / GA / AI / TravelTime};

  \draw[arr] (domain) -- (sched);
  \draw[arr] (sched) -- (app);
  \draw[arr] (proto) -- (domain);
  \draw[arr, dashed] (app.south) to[bend left=20] node[right, font=\scriptsize]{adapters} (proto.east);
\end{tikzpicture}
\caption{Dipendenze Fase 8: package puro + adapter nell'app.}
\end{figure}

\section{Contenuto package}

Symlink da \repo{Services/Scheduling/} (7 file):

\begin{itemize}
  \item \texttt{PathResolver}, \texttt{KinematicCalculator}
  \item \texttt{TaktEngine}, \texttt{TaktEngineScheduleGeneration}
  \item \texttt{ScheduleRefresher}, \texttt{VehicleSuitabilityEngine}
  \item \texttt{ScheduleOptimizationPipeline}
\end{itemize}

\section{Protocolli (FDCDomain)}

\begin{table}[htbp]
\centering
\caption{Boundary tra package e app.}
\begin{tabularx}{\textwidth}{lX}
  \toprule
  \textbf{Protocollo} & \textbf{Implementazione app} \\
  \midrule
  \texttt{ScheduleConflictDetecting} & \texttt{ConflictManagerAdapter} \\
  \texttt{ScheduleGeneticOptimizing} & \texttt{GeneticOptimizerAdapter} \\
  \texttt{ScheduleAIOptimizing} & \texttt{ScheduleAIResolverAdapter} \\
  \texttt{TrainTravelTimeCalculating} & \texttt{FDCSchedulerTravelTimeAdapter} \\
  \bottomrule
\end{tabularx}
\end{table}

\section{Shell app}

Restano nel target app:

\begin{itemize}
  \item \texttt{ScheduleGenerationEngine} --- orchestratore UI
  \item \texttt{ScheduleAIResolver} --- RailwayAIService
  \item \texttt{AppSchedulingAdapters.swift} --- conformità protocolli
  \item \texttt{SchedulingBridge.swift} --- \texttt{@\_exported import FDCScheduling}
\end{itemize}

\section{Test}

\begin{itemize}
  \item \textbf{6 test} \texttt{FDCDomainTests}
  \item \textbf{4 test} \texttt{FDCSchedulingTests}
  \item \textbf{24 test} app (invariati, stub travel-time nei test)
\end{itemize}

\section{Nota tecnica}

I path con carattere \texttt{+} nelle \texttt{membershipExceptions} del
\texttt{pbxproj} corrompono il parser Xcode; i file sono stati rinominati
(es.\ \texttt{DateScheduling.swift}, \texttt{TaktEngineScheduleGeneration.swift}).


\part{Moduli applicativi}
\chapter{Visualizzazione infrastruttura hub/AV}
\label{ch:infrastructure}

\section{Modello hub}

Un \textbf{hub} unisce due nodi fisici distinti ma logicamente la stessa
stazione ferroviaria:

\begin{itemize}
  \item \textbf{Stazione classica} --- nodo \emph{parent} al centro del
        quadrato schematico (\texttt{parentHubId == nil}).
  \item \textbf{Stazione AV} --- nodo \emph{satellite} spostato su uno
        degli 8 punti del quadrato (4 vertici + 4 metà lati), tramite
        \texttt{hubOffsetDirection}.
\end{itemize}

Le coordinate geografiche del satellite ereditano quelle del parent
(calcolo distanze in \texttt{NetworkTopologyService}); la mappa applica
l'offset canvas in \texttt{MapGeometryEngine.finalPosition}.

\section{Regole di routing binari}

\texttt{HubTopology} (Domain) risolve gli endpoint fisici in base al
tipo traccia:

\begin{center}
\begin{tabular}{lll}
  \toprule
  \textbf{Tipo binario} & \textbf{Endpoint} & \textbf{Motivo} \\
  \midrule
  \texttt{regional}, \texttt{single} & parent (centro) & binario classico \\
  \texttt{highSpeed} & satellite AV & binario AV separato \\
  \bottomrule
\end{tabular}
\end{center}

La risoluzione avviene in \texttt{RailroadNetwork.addEdge} prima della
persistenza: l'utente può selezionare la stazione logica; il sistema
collega il nodo fisico corretto.

\section{Rendering}

\begin{itemize}
  \item \textbf{Pallini rossi} --- entrambi i nodi hub
        (\texttt{drawHubStationDot}).
  \item \textbf{Linea a 8} --- curva chiusa tra classico e satellite
        (\texttt{drawHubPair}).
  \item \textbf{Etichetta} --- nome del parent sotto il gruppo.
  \item I binari AV usano le coordinate del satellite, quindi non si
        sovrappongono ai binari classici al centro.
\end{itemize}

\section{Workflow editor}

Dall'inspector stazione (\texttt{StationInlineEditor}):
\begin{enumerate}
  \item Selezionare la stazione classica (parent).
  \item Pulsante \emph{Crea stazione AV satellite}.
  \item Regolare \texttt{hubOffsetDirection} sul satellite.
  \item Creare binari AV tra satelliti; binari regionali tra parent.
\end{enumerate}

\section{Sorgenti}

\begin{itemize}
  \item \repo{Domain/Services/HubTopology.swift}
  \item \repo{Domain/Model/Node.swift} --- \texttt{HubOffsetDirection}
  \item \repo{Rendering/RailwayRenderer.swift}
  \item \repo{Rendering/MapGeometryEngine.swift}
  \item \repo{UI/Views/RailwayMap/InfrastructureCanvas.swift}
\end{itemize}


% Prossimi capitoli (aggiungere man mano):
% \input{modules/05-editor}
% \input{modules/06-simulator}
% \input{modules/07-railway-ai}

\appendix
\chapter{Indice dei sorgenti}
\begin{longtable}{ll}
  \toprule
  \textbf{Percorso} & \textbf{Modulo} \\
  \midrule
  \endhead
  \repo{Domain/Services/}     & NetworkTopologyService \\
  \repo{Domain/Pathfinding/}  & NetworkPathfinder \\
  \repo{Domain/Model/}        & FDCDomain \\
  \repo{Services/Scheduling/} & FDCScheduling \\
  \repo{Editor/}             & Editor \textit{(da documentare)} \\
  \repo{Scheduler/}          & UI Scheduler \\
  \repo{Simulator/}          & Simulator \textit{(da documentare)} \\
  \repo{RailwayAIService/}   & AI \textit{(da documentare)} \\
  \repo{Infrastructure/}     & I/O \textit{(da documentare)} \\
  \repo{Packages/FDCDomain/} & SPM Domain \\
  \repo{Packages/FDCScheduling/} & SPM Scheduling \\
  \bottomrule
\end{longtable}

\end{document}
